\section{概率空间、条件概率与随机变量}


这一讲从最基础的概率语言开始：先把“样本空间—事件—概率”三件事说清楚，再引入条件概率、条件分布、独立性，最后把随机变量定义为样本空间上的函数，并给出 Bernoulli 与二项分布。

\textbf{本节用到的记号：} $\Omega$ 为样本空间，$\omega$ 为样本点；$A,B\subseteq\Omega$ 为事件；$P$ 为概率测度；$X:\Omega\to S$ 为随机变量；$p_X(x)=P(X=x)$ 为分布；$\mathbf 1_A$ 为事件 $A$ 的指示变量。

\subsection{离散概率空间}


本课程主要考虑\textbf{离散概率}：样本空间 $\Omega$ 至多可数，每个样本点 $\omega$ 有一个非负权重 $p(\omega)$，满足

\[
\sum_{\omega\in\Omega}p(\omega)=1.
\]

事件 $A\subseteq\Omega$ 的概率定义为

\[
P(A)=\sum_{\omega\in A}p(\omega),\qquad P(\varnothing)=0,\quad P(\Omega)=1.
\]

这就是“把事件拆成不可再分的样本点，再求和”。

\textbf{基本性质.}

\begin{itemize}
\item 补集：$P(A^c)=1-P(A)$。
\item 单调性：$A\subseteq B\Rightarrow P(A)\le P(B)$。
\item 加法：若 $A\cap B=\varnothing$，则 $P(A\cup B)=P(A)+P(B)$。
\item 容斥：$P(A\cup B)=P(A)+P(B)-P(A\cap B)$。
\item 次可加（union bound）：$P(A\cup B)\le P(A)+P(B)$。
\end{itemize}


\textbf{独立重复试验的样本空间.} 若一次试验有 $m$ 个等可能结果，独立做 $n$ 次，则样本空间大小为 $m^n$。抛 $n$ 枚硬币、掷 $n$ 次骰子都属于这种情况。

\subsection{条件概率与条件分布}


当已知 $B$ 发生，样本空间从 $\Omega$ 缩小到 $B$。定义

\[
P(A\mid B)=\frac{P(A\cap B)}{P(B)}\qquad(P(B)>0).
\]

由此得到乘法公式

\[
P(A\cap B)=P(B)P(A\mid B)=P(A)P(B\mid A).
\]

\textbf{条件分布.} 固定事件 $B$ 后，可以重新归一化，得到一个新的概率分布：

\[
p_B(\omega)=
\begin{cases}
\dfrac{p(\omega)}{P(B)}, & \omega\in B,\\[4pt]
0, & \omega\notin B.
\end{cases}
\]

它满足 $\sum_{\omega}p_B(\omega)=1$，所以“以 $B$ 为条件”可以把原来的概率空间换成 $(B,p_B)$。

\textbf{全概率公式.} 若 $B_1,\dots,B_n$ 是 $\Omega$ 的一个划分（两两不交、并为 $\Omega$），且每个 $P(B_i)>0$，则

\[
P(A)=\sum_{i=1}^{n}P(B_i)P(A\mid B_i).
\]

\textbf{Bayes 公式.} 在同样条件下，

\[
P(B_i\mid A)
=
\frac{P(B_i)P(A\mid B_i)}
{\sum_{j=1}^{n}P(B_j)P(A\mid B_j)}.
\]

分子是“先到 $B_i$、再到 $A$”，分母是把所有可能的“先到 $B_j$”都加起来。

\subsection{独立性}


两个事件 $A,B$ 独立，定义为

\[
P(A\cap B)=P(A)P(B).
\]

当 $P(B)>0$ 时，这也等价于 $P(A\mid B)=P(A)$：知道 $B$ 发生不改变 $A$ 的概率。

多个事件 $A_1,\dots,A_n$ 相互独立，要求对任意子集 $S\subseteq[n]$ 都有

\[
P\Bigl(\bigcap_{i\in S}A_i\Bigr)=\prod_{i\in S}P(A_i).
\]

注意：\textbf{两两独立不等于相互独立}；三个事件还需要满足三三交集的条件。

\subsection{随机变量与分布}


随机变量是样本空间上的函数

\[
X:\Omega\to S,
\]

它把每个样本点 $\omega$ 映射成一个值 $X(\omega)$。离散情形下，$X$ 的分布为

\[
p_X(x)=P(X=x)=P(\{\omega:X(\omega)=x\}).
\]

记 $X\sim p_X$。\textbf{随机变量的分布}才是我们真正关心的信息；样本空间只是“背后的概率来源”。

\textbf{指示变量.} 对事件 $A$，

\[
\mathbf 1_A(\omega)=
\begin{cases}
1, & \omega\in A,\\
0, & \omega\notin A.
\end{cases}
\]

它是随机变量，且

\[
E[\mathbf 1_A]=P(A).
\]

\textbf{随机变量独立.} $X,Y$ 独立，当且仅当对任意 $x,y$，

\[
P(X=x,Y=y)=P(X=x)P(Y=y).
\]

等价地，联合分布等于边缘分布的乘积。多个随机变量相互独立也要求任意子族的联合分布等于各自边缘分布之积。

\subsection{Bernoulli 与二项分布}


\textbf{Bernoulli 分布.} $X\sim\mathrm{Bern}(p)$，取值 $0,1$：

\[
P(X=1)=p,\qquad P(X=0)=1-p.
\]

\textbf{二项分布.} 若 $X_1,\dots,X_n$ 独立同分布，$X_i\sim\mathrm{Bern}(p)$，令

\[
X=X_1+\cdots+X_n,
\]

则 $X$ 是“$n$ 次独立试验中成功的次数”，

\[
P(X=k)=\binom nk p^k(1-p)^{n-k},\qquad k=0,1,\dots,n.
\]

记作 $X\sim\mathrm{Bin}(n,p)$。

\textbf{直观理解.} 先选出哪 $k$ 次成功，共 $\binom nk$ 种；每一种指定成功位置的概率都是 $p^k(1-p)^{n-k}$。

\textbf{例.} 抛 $10$ 枚公平硬币，正面数 $X\sim\mathrm{Bin}(10,1/2)$，则

\[
P(X=3)=\binom{10}{3}\left(\frac12\right)^{10}=\frac{120}{1024}.
\]

\textbf{例.} 掷 $n$ 次骰子，统计出现 $6$ 的次数，得到 $\mathrm{Bin}(n,1/6)$。

\subsection{均匀分布}


若 $X$ 在有限集合 $S$ 上等可能取值，则

\[
P(X=x)=
\begin{cases}
\dfrac{1}{|S|}, & x\in S,\\[4pt]
0, & x\notin S.
\end{cases}
\]

这就是\textbf{离散均匀分布}。掷一枚公平骰子，点数 $X$ 在 $\{1,\dots,6\}$ 上均匀分布。

若 $S=\{1,2,\dots,m\}$，则

\[
E[X]=\frac{m+1}{2},\qquad
\operatorname{Var}(X)=\frac{m^2-1}{12}.
\]

均匀分布是“没有任何取值偏好”的最基本模型；后面很多“等可能样本点”的计数都默认均匀分布。
